\documentclass[aspectratio=169,10pt]{beamer}

\usepackage[T1]{fontenc}

\usepackage[utf8]{inputenc}

\usepackage[spanish,mexico]{babel}

\usepackage{lmodern}

\usepackage{microtype}

\usepackage{tikz}

\usepackage{graphicx}

\usepackage{booktabs}

\usepackage{array}

\usepackage{tabularx}

\usepackage{ragged2e}

\usetikzlibrary{positioning,arrows.meta,calc,shapes.geometric,fit,backgrounds}

\graphicspath{{eld_figs/}}



% Visual system adapted from the reference deck

\definecolor{CambridgeNavy}{HTML}{18334B}

\definecolor{JudgeTeal}{HTML}{267D83}

\definecolor{CambridgeGold}{HTML}{A68449}

\definecolor{Quiet}{HTML}{63707B}

\definecolor{Pale}{HTML}{F3F6F6}

\definecolor{IdeaPale}{HTML}{EDF5F4}

\definecolor{IdeaInk}{HTML}{244653}

\definecolor{SafeRed}{HTML}{A33A3A}

\definecolor{ExecBlue}{HTML}{246FA8}

\definecolor{SoftAmber}{HTML}{FFF7ED}

\definecolor{SoftRed}{HTML}{FDF0F0}

\definecolor{SoftBlue}{HTML}{EAF2F8}



\setbeamercolor{normal text}{fg=IdeaInk,bg=white}

\setbeamercolor{frametitle}{fg=CambridgeNavy,bg=white}

\setbeamercolor{structure}{fg=JudgeTeal}

\setbeamercolor{alerted text}{fg=CambridgeGold}

\setbeamertemplate{navigation symbols}{}

\setbeamertemplate{itemize item}{\color{JudgeTeal}$\blacktriangleright$}

\setbeamertemplate{itemize subitem}{\color{CambridgeGold}$\bullet$}

\setbeamertemplate{frametitle}{%

  \vspace{0.18cm}\insertframetitle\par\vspace{-0.11cm}%

  {\color{CambridgeGold}\rule{1.35cm}{0.7pt}}}

\setbeamertemplate{footline}{%

  \leavevmode\hbox{%

  \begin{beamercolorbox}[wd=\paperwidth,ht=2.45ex,dp=1.05ex,right]{normal text}%

  \color{Quiet}\scriptsize Escuela Libre de Derecho \enspace\textcolor{CambridgeGold}{|}\enspace Del conocimiento jurídico a las instituciones inteligentes

  \hfill\insertframenumber/18\hspace{0.38cm}

  \end{beamercolorbox}}}



\newcommand{\key}[1]{%
  \begin{beamercolorbox}[rounded=true,sep=0.62em,wd=\textwidth]{block body}%
  \color{JudgeTeal}\bfseries #1
  \end{beamercolorbox}}

\newcommand{\pill}[1]{%
  \tikz[baseline=-0.6ex]\node[rounded corners=2pt,fill=IdeaPale,text=JudgeTeal,
  inner xsep=5pt,inner ysep=2.5pt,font=\scriptsize\bfseries]{#1};}

\newcommand{\stagebar}[2]{%
  \vspace{-0.08cm}
  \begin{tikzpicture}
    \node[rounded corners=1.5pt,fill=CambridgeNavy,text=white,
      inner xsep=5pt,inner ysep=2.2pt,font=\tiny\bfseries] {#1};
    \node[anchor=west,text=Quiet,font=\tiny\bfseries] at (1.52,0) {#2};
    \draw[Quiet!28,line width=.45pt] (4.15,0)--(13.05,0);
  \end{tikzpicture}\par\vspace{0.04cm}}

\newcommand{\sidecard}[3]{%
  \begin{beamercolorbox}[rounded=true,sep=0.50em,wd=3.28cm]{block body}
    {\color{#1}\scriptsize\bfseries #2}\par\smallskip
    {\color{IdeaInk}\scriptsize #3}
  \end{beamercolorbox}}

\newcommand{\sectioncard}[3]{%
  \begin{beamercolorbox}[rounded=true,sep=0.38em,wd=3.18cm]{block body}
    {\color{#1}\tiny\bfseries #2}\par\smallskip
    {\color{IdeaInk}\tiny #3}
  \end{beamercolorbox}}

\newcommand{\minicard}[3]{%
  \begin{tikzpicture}
    \node[draw=#1!55,rounded corners=3pt,fill=#1!8,text width=3.42cm,minimum height=2.20cm,inner sep=7pt,align=left] {
      {\color{#1}\bfseries\small #2}\par\smallskip
      {\scriptsize #3}};
  \end{tikzpicture}}

\newcommand{\sectionslide}[7]{%
\begin{frame}[t]{#2}
  \stagebar{#1}{#3}
  \begin{columns}[T,onlytextwidth]
    \column{.745\textwidth}
      \begin{tikzpicture}
        \node[draw=CambridgeNavy!30,line width=.65pt,fill=white,
          minimum width=10.18cm,minimum height=5.72cm,
          inner sep=0pt,outer sep=0pt] (canvas) {};
        \node[draw=JudgeTeal!65,rounded corners=3pt,fill=IdeaPale,text width=8.45cm,minimum height=1.24cm,inner sep=8pt,align=left] at (0,1.52)
          {\color{JudgeTeal}\bfseries\small LECTURA\par\color{IdeaInk}\scriptsize #5};
        \node[draw=CambridgeGold!70,rounded corners=3pt,fill=SoftAmber,text width=8.45cm,minimum height=1.24cm,inner sep=8pt,align=left] at (0,0)
          {\color{CambridgeGold}\bfseries\small TESIS\par\color{IdeaInk}\scriptsize #6};
        \node[draw=SafeRed!60,rounded corners=3pt,fill=SoftRed,text width=8.45cm,minimum height=1.24cm,inner sep=8pt,align=left] at (0,-1.52)
          {\color{SafeRed}\bfseries\small CUIDADO\par\color{IdeaInk}\scriptsize #7};
        \draw[-{Latex},thick,color=Quiet!70] (-3.5,0.86)--(-3.5,0.35);
        \draw[-{Latex},thick,color=Quiet!70] (-3.5,-0.66)--(-3.5,-1.17);
      \end{tikzpicture}
    \column{.235\textwidth}
      \sectioncard{JudgeTeal}{LECTURA}{#5}
      \vspace{0.10cm}
      \sectioncard{CambridgeGold}{TESIS}{#6}
      \vspace{0.10cm}
      \sectioncard{SafeRed}{CUIDADO}{#7}
  \end{columns}
\end{frame}}

\begin{document}



% 1

{

\setbeamercolor{background canvas}{bg=CambridgeNavy}

\begin{frame}[plain]

\begin{tikzpicture}[remember picture,overlay]

  \fill[JudgeTeal] (current page.south west) rectangle ([xshift=0.16cm]current page.north west);

  \fill[JudgeTeal!15!CambridgeNavy] ([xshift=-4.35cm]current page.south east) rectangle (current page.north east);

  \draw[CambridgeGold,line width=.9pt] ([xshift=1.0cm,yshift=-2.05cm]current page.north west) -- ++(1.65cm,0);

\end{tikzpicture}

\vspace{0.58cm}

{\color{JudgeTeal}\scriptsize\bfseries ESCUELA LIBRE DE DERECHO \enspace·\enspace INTELIGENCIA ARTIFICIAL Y DERECHO}\par

\vspace{0.72cm}

\begin{minipage}{11.35cm}

{\color{white}\fontsize{20}{23}\selectfont\bfseries

Del conocimiento jurídico\\[0.04cm]a las instituciones inteligentes}

\end{minipage}\par

\vspace{0.34cm}

{\color{CambridgeGold}\fontsize{12}{15}\selectfont El Derecho frente a una nueva forma de inteligencia}\par

\vspace{0.74cm}

\begin{minipage}{9.75cm}

{\color{white!88}\normalsize La pregunta no es sólo qué puede hacer una máquina, sino bajo qué fuentes, límites, autoridad y responsabilidad puede hacerlo.}

\end{minipage}

\vfill

{\color{white!82}\scriptsize\bfseries Alejandro Reynoso}\par

\vspace{0.08cm}

{\color{white!62}\tiny Judge Business School · University of Cambridge}

\vspace{0.38cm}

\end{frame}

}



% 2

\begin{frame}{Una tesis en seis movimientos}

\key{La conferencia no trata de usar ChatGPT para escribir más rápido. Trata de cómo el Derecho gobierna una nueva forma de inteligencia que entra en sus procesos.}

\vspace{0.25cm}

\begin{center}

\begin{tikzpicture}[node distance=.17cm,every node/.style={align=left}]

  \node[draw=CambridgeNavy,rounded corners=3pt,fill=Pale,minimum width=2.05cm,minimum height=1.75cm,text width=1.72cm,inner sep=6pt] (a) {{\color{CambridgeNavy}\bfseries\small 01}\par\scriptsize Conocimiento jurídico};

  \node[draw=ExecBlue,rounded corners=3pt,fill=SoftBlue,minimum width=2.05cm,minimum height=1.75cm,text width=1.72cm,inner sep=6pt,right=of a] (b) {{\color{ExecBlue}\bfseries\small 02}\par\scriptsize Memoria};

  \node[draw=JudgeTeal,rounded corners=3pt,fill=IdeaPale,minimum width=2.05cm,minimum height=1.75cm,text width=1.72cm,inner sep=6pt,right=of b] (c) {{\color{JudgeTeal}\bfseries\small 03}\par\scriptsize Integración};

  \node[draw=CambridgeGold,rounded corners=3pt,fill=SoftAmber,minimum width=2.05cm,minimum height=1.75cm,text width=1.72cm,inner sep=6pt,right=of c] (d) {{\color{CambridgeGold}\bfseries\small 04}\par\scriptsize Sistema};

  \node[draw=SafeRed,rounded corners=3pt,fill=SoftRed,minimum width=2.05cm,minimum height=1.75cm,text width=1.72cm,inner sep=6pt,right=of d] (e) {{\color{SafeRed}\bfseries\small 05}\par\scriptsize Autoridad};

  \node[draw=Quiet,rounded corners=3pt,fill=Pale,minimum width=2.05cm,minimum height=1.75cm,text width=1.72cm,inner sep=6pt,right=of e] (f) {{\color{Quiet}\bfseries\small 06}\par\scriptsize Gobierno};

  \foreach \x/\y in {a/b,b/c,c/d,d/e,e/f}{\draw[-{Latex},thick,color=Quiet] (\x)--(\y);}  

\end{tikzpicture}

\end{center}

\vspace{0.30cm}

\begin{columns}[T,onlytextwidth]

\column{.31\textwidth}\sidecard{JudgeTeal}{PUNTO DE PARTIDA}{Una pregunta a una máquina produce una respuesta convincente.}

\column{.31\textwidth}\sidecard{CambridgeGold}{GIRO CONCEPTUAL}{La respuesta obliga a diseñar memoria, integración, límites y trazabilidad.}

\column{.31\textwidth}\sidecard{SafeRed}{PUNTO DE LLEGADA}{Una institución debe gobernar inteligencia, no sólo consumirla.}

\end{columns}

\end{frame}



% 3

\begin{frame}{El problema ya ocurrió: Mata v. Avianca}

\stagebar{APERTURA}{De una cita inexistente a una pregunta institucional}

\begin{columns}[T,onlytextwidth]

\column{.58\textwidth}

\begin{tikzpicture}

  \node[draw=CambridgeNavy!30,rounded corners=3pt,fill=Pale,text width=7.20cm,inner sep=10pt] (case) {

    {\color{CambridgeNavy}\Large\bfseries En 2023, unas sentencias inexistentes llegaron a un tribunal.}\par\medskip

    \small En \textit{Mata v. Avianca}, los abogados presentaron citas generadas por ChatGPT. El tribunal sancionó la falta de verificación: la tecnología era nueva; la responsabilidad profesional no había desaparecido.};

  \node[rounded corners=2pt,fill=SoftRed,draw=SafeRed!55,text width=6.65cm,inner sep=8pt,anchor=north west] at ([yshift=-.35cm]case.south west) {

  \color{SafeRed}\bfseries\small La pregunta central\par\smallskip

  \color{IdeaInk}\scriptsize ¿Qué debe existir alrededor de la inteligencia artificial para que pueda usarse responsablemente en Derecho?};

\end{tikzpicture}

\column{.38\textwidth}

\sidecard{ExecBlue}{DÓNDE ESTAMOS}{AI Act vigente; adopción acelerada; políticas internas rezagadas.}

\vspace{0.18cm}

\sidecard{CambridgeGold}{FOTOGRAFÍA DEL MOMENTO}{La capacidad avanza más rápido que su gobernanza.}

\vspace{0.18cm}

\sidecard{SafeRed}{RIESGO INMEDIATO}{Una respuesta convincente puede carecer de fuente o autoridad.}

\end{columns}

\vspace{0.15cm}

\centering\pill{Capacidad no es autoridad}\hspace{0.18cm}\pill{Precisión no es legitimidad}\hspace{0.18cm}\pill{Verificar sigue siendo deber profesional}

\end{frame}



% 4--9 section anchor images

\sectionslide{SECCIÓN 1}{De un caso real a una pregunta ineludible}{Qué pasó · dónde estamos hoy · qué está en juego}{sec_1.png}{El riesgo ya ocurrió: una respuesta falsa llegó a un tribunal.}{La IA exige fuentes, verificación y responsabilidad.}{La forma correcta no prueba autoridad.}



\sectionslide{SECCIÓN 2}{De la conversación al segundo cerebro}{Del chat individual a la memoria institucional}{sec_2.png}{Del chat individual a memoria con relaciones y procedencia.}{Memoria útil = vínculos, procedencia y control.}{No todo expediente puede entrar al sistema.}



\sectionslide{SECCIÓN 3}{Integrar lo que la especialización fragmentó}{De silos jurídicos a visión de conjunto}{sec_3.png}{Las ramas ayudan a estudiar; los problemas cruzan fronteras.}{Recuperar integración sin eliminar profundidad.}{Derecho comparado exige función y contexto.}



\sectionslide{SECCIÓN 4}{Del chat al sistema}{Herramientas · agentes · arneses}{sec_4.png}{Modelos + datos + reglas + memoria = sistema.}{La inteligencia reside en la arquitectura.}{Más capacidad exige más arnés.}



\sectionslide{SECCIÓN 5}{Capacidad no es autoridad}{Facultades · responsabilidad · límites}{sec_5.png}{Una máquina puede hacer; no por ello puede decidir.}{Precisión técnica no es legitimidad institucional.}{La autoridad viene del Derecho.}



\sectionslide{SECCIÓN 6}{Cuando también cambia la organización}{Constelaciones · adaptación · gobierno del cambio}{sec_6.png}{El sistema también puede proponer cambios al proceso.}{No evitar el cambio: gobernarlo.}{Separar propuesta, evaluación y autorización.}



% 10

\begin{frame}{Caso I: enjambre de 50 agentes para la LMV}

\stagebar{LABORATORIO}{De discurso conceptual a arquitectura operativa}

\key{Tarea: identificar áreas de mejora de la Ley del Mercado de Valores y producir un anteproyecto limitado a 20 artículos reformados, adicionados o derogados.}

\vspace{0.25cm}

\begin{center}

\begin{tikzpicture}[node distance=.25cm,every node/.style={align=center}]

  \node[draw=CambridgeNavy,rounded corners=3pt,fill=Pale,minimum width=2.4cm,minimum height=1.1cm,font=\scriptsize\bfseries] (corpus) {LMV oficial\\+ corpus autorizado};

  \node[draw=JudgeTeal,rounded corners=3pt,fill=IdeaPale,minimum width=2.4cm,minimum height=1.1cm,font=\scriptsize\bfseries,right=of corpus] (agents) {50 agentes\\8 familias};

  \node[draw=ExecBlue,rounded corners=3pt,fill=SoftBlue,minimum width=2.4cm,minimum height=1.1cm,font=\scriptsize\bfseries,right=of agents] (const) {Constelaciones\\por problema};

  \node[draw=CambridgeGold,rounded corners=3pt,fill=SoftAmber,minimum width=2.4cm,minimum height=1.1cm,font=\scriptsize\bfseries,right=of const] (crit) {Crítica\\adversarial};

  \node[draw=SafeRed,rounded corners=3pt,fill=SoftRed,minimum width=2.4cm,minimum height=1.1cm,font=\scriptsize\bfseries,right=of crit] (bill) {20 artículos\\anteproyecto};

  \foreach \x/\y in {corpus/agents,agents/const,const/crit,crit/bill}{\draw[-{Latex},thick,color=Quiet] (\x)--(\y);}  

\end{tikzpicture}

\end{center}

\vspace{0.24cm}

\begin{columns}[T,onlytextwidth]

\column{.31\textwidth}\sidecard{JudgeTeal}{NO ES UN PROMPT}{El corpus define la frontera de conocimiento y la procedencia de cada artículo.}

\column{.31\textwidth}\sidecard{CambridgeGold}{NO ES VOTACIÓN}{El sistema conserva desacuerdos, objeciones y sesgos institucionales explícitos.}

\column{.31\textwidth}\sidecard{SafeRed}{NO ES AUTORIDAD}{La salida es un anteproyecto para revisión humana, no una iniciativa autorizada.}

\end{columns}

\end{frame}



% 11

\begin{frame}{El enjambre: perspectivas con sesgos declarados}

\stagebar{DISEÑO}{Diversidad institucional sin fingir neutralidad}

\begin{columns}[T,onlytextwidth]

\column{.52\textwidth}

\begin{tikzpicture}[node distance=.10cm]
\node[draw=JudgeTeal!65,fill=JudgeTeal!8,rounded corners=2pt,text width=6.12cm,minimum width=6.50cm,minimum height=.50cm,anchor=west,align=left,inner sep=3.5pt] at (0,0) {\tiny\makebox[1.48cm][l]{\color{JudgeTeal}\bfseries Abogados}\color{IdeaInk}práctica · técnica · ejecutabilidad};
\node[draw=CambridgeGold!65,fill=CambridgeGold!8,rounded corners=2pt,text width=6.12cm,minimum width=6.50cm,minimum height=.50cm,anchor=west,align=left,inner sep=3.5pt] at (0,-0.67) {\tiny\makebox[1.48cm][l]{\color{CambridgeGold}\bfseries Legisladores}\color{IdeaInk}política pública · constitucionalidad};
\node[draw=SafeRed!65,fill=SafeRed!8,rounded corners=2pt,text width=6.12cm,minimum width=6.50cm,minimum height=.50cm,anchor=west,align=left,inner sep=3.5pt] at (0,-1.34) {\tiny\makebox[1.48cm][l]{\color{SafeRed}\bfseries Reguladores}\color{IdeaInk}supervisión · integridad de mercado};
\node[draw=ExecBlue!65,fill=ExecBlue!8,rounded corners=2pt,text width=6.12cm,minimum width=6.50cm,minimum height=.50cm,anchor=west,align=left,inner sep=3.5pt] at (0,-2.01) {\tiny\makebox[1.48cm][l]{\color{ExecBlue}\bfseries Emisoras}\color{IdeaInk}acceso · costos · proporcionalidad};
\node[draw=CambridgeNavy!65,fill=CambridgeNavy!8,rounded corners=2pt,text width=6.12cm,minimum width=6.50cm,minimum height=.50cm,anchor=west,align=left,inner sep=3.5pt] at (0,-2.68) {\tiny\makebox[1.48cm][l]{\color{CambridgeNavy}\bfseries Bolsas}\color{IdeaInk}listado · liquidez · infraestructura};
\node[draw=Quiet!65,fill=Quiet!8,rounded corners=2pt,text width=6.12cm,minimum width=6.50cm,minimum height=.50cm,anchor=west,align=left,inner sep=3.5pt] at (0,-3.35) {\tiny\makebox[1.48cm][l]{\color{Quiet}\bfseries Intermediarios}\color{IdeaInk}colocación · clientes · responsabilidad};
\node[draw=JudgeTeal!65,fill=JudgeTeal!8,rounded corners=2pt,text width=6.12cm,minimum width=6.50cm,minimum height=.50cm,anchor=west,align=left,inner sep=3.5pt] at (0,-4.02) {\tiny\makebox[1.48cm][l]{\color{JudgeTeal}\bfseries Inversionistas}\color{IdeaInk}revelación · gobierno · minorías};
\node[draw=CambridgeGold!65,fill=CambridgeGold!8,rounded corners=2pt,text width=6.12cm,minimum width=6.50cm,minimum height=.50cm,anchor=west,align=left,inner sep=3.5pt] at (0,-4.69) {\tiny\makebox[1.48cm][l]{\color{CambridgeGold}\bfseries Comparado}\color{IdeaInk}experiencia internacional · academia};
\end{tikzpicture}

\column{.43\textwidth}

\sidecard{JudgeTeal}{PRINCIPIO}{Un buen enjambre no oculta intereses; los hace visibles para tratarlos como información del proceso.}

\vspace{0.15cm}

\sidecard{CambridgeGold}{MÉTODO}{Diagnóstico independiente → constelaciones mixtas → propuesta → crítica externa.}

\vspace{0.15cm}

\sidecard{SafeRed}{CONTROL}{El consenso artificial no sustituye evidencia, técnica legislativa ni revisión humana.}

\end{columns}

\end{frame}



% 12

\begin{frame}{Arquitectura de gobierno: inner loop y outer loop}

\stagebar{ARNÉS}{El sistema delibera, pero no se autoautoriza}

\begin{center}

\begin{tikzpicture}[node distance=.20cm,every node/.style={align=center}]

\node[draw=CambridgeNavy,rounded corners=3pt,fill=Pale,minimum width=2.35cm,minimum height=.92cm,font=\scriptsize\bfseries] (a) {Diagnóstico};

\node[draw=JudgeTeal,rounded corners=3pt,fill=IdeaPale,minimum width=2.35cm,minimum height=.92cm,font=\scriptsize\bfseries,right=of a] (b) {Propuesta};

\node[draw=CambridgeGold,rounded corners=3pt,fill=SoftAmber,minimum width=2.35cm,minimum height=.92cm,font=\scriptsize\bfseries,right=of b] (c) {Crítica};

\node[draw=ExecBlue,rounded corners=3pt,fill=SoftBlue,minimum width=2.35cm,minimum height=.92cm,font=\scriptsize\bfseries,right=of c] (d) {Selección};

\node[draw=SafeRed,rounded corners=3pt,fill=SoftRed,minimum width=2.35cm,minimum height=.92cm,font=\scriptsize\bfseries,right=of d] (e) {Redacción};

\foreach \x/\y in {a/b,b/c,c/d,d/e}{\draw[-{Latex},thick,color=Quiet] (\x)--(\y);}  

\draw[JudgeTeal,thick,-{Latex}] ([yshift=-.45cm]e.south) -- ++(0,-.55) -| ([yshift=-.45cm]a.south);

\begin{scope}[on background layer]

\node[draw=CambridgeNavy!45,rounded corners=4pt,fill=white,fit=(a)(b)(c)(d)(e),inner sep=8pt,label={[font=\scriptsize\bfseries,text=CambridgeNavy]above:INNER LOOP}]{};

\end{scope}

\node[draw=SafeRed!70,rounded corners=3pt,fill=SoftRed,minimum width=8.15cm,minimum height=.78cm,anchor=north] (outer) at ([yshift=-1.60cm]c.south) {\scriptsize\bfseries OUTER LOOP: cobertura · concentración de influencia · trazabilidad · validación · objeciones pendientes};

\draw[-{Latex},thick,color=SafeRed] (outer.north) -- ([yshift=-.26cm]c.south);

\end{tikzpicture}

\end{center}

\vspace{0.25cm}

\begin{columns}[T,onlytextwidth]

\column{.31\textwidth}\sidecard{JudgeTeal}{PROPONE}{Agentes y constelaciones generan diagnósticos y alternativas.}

\column{.31\textwidth}\sidecard{CambridgeGold}{EVALÚA}{Crítica adversarial y outer loop identifican fallas.}

\column{.31\textwidth}\sidecard{SafeRed}{AUTORIZA}{La selección de 20 artículos requiere puerta humana explícita.}

\end{columns}

\end{frame}



% 13

\begin{frame}{Producto esperado: 20 artículos y un expediente de trazabilidad}

\stagebar{SALIDA}{No sólo texto legislativo: genealogía de cada cambio}

\key{La salida no es “un texto bonito”: es un paquete revisable que conecta cada artículo con evidencia, objeciones, responsables y validación.}

\vspace{0.22cm}

\begin{center}

\begin{tikzpicture}[node distance=.22cm,every node/.style={align=left}]

\node[draw=CambridgeNavy,rounded corners=3pt,fill=Pale,text width=3.90cm,minimum height=2.0cm,inner sep=8pt] (a) {\color{CambridgeNavy}\bfseries Anteproyecto\par\smallskip\color{IdeaInk}\scriptsize 20 artículos con operación legislativa, problema, evidencia y redacción propuesta.};

\node[draw=JudgeTeal,rounded corners=3pt,fill=IdeaPale,text width=3.90cm,minimum height=2.0cm,inner sep=8pt,right=of a] (b) {\color{JudgeTeal}\bfseries Validación\par\smallskip\color{IdeaInk}\scriptsize Artículos existentes verificados; adiciones marcadas; objeciones visibles.};

\node[draw=CambridgeGold,rounded corners=3pt,fill=SoftAmber,text width=3.90cm,minimum height=2.0cm,inner sep=8pt,right=of b] (c) {\color{CambridgeGold}\bfseries Trazabilidad\par\smallskip\color{IdeaInk}\scriptsize Quién propuso, quién criticó, qué fuente sostuvo y qué quedó pendiente.};

\foreach \x/\y in {a/b,b/c}{\draw[-{Latex},thick,color=Quiet] (\x)--(\y);}

\end{tikzpicture}

\end{center}

\vspace{0.30cm}

\begin{columns}[T,onlytextwidth]

\column{.31\textwidth}\sidecard{ExecBlue}{MODELO}{Claude Haiku 4.5 para todos los agentes; diversidad por rol y mandato.}

\column{.31\textwidth}\sidecard{CambridgeGold}{CORPUS}{La LMV oficial manda; el comparado genera hipótesis.}

\column{.31\textwidth}\sidecard{SafeRed}{LEYENDA}{Anteproyecto para revisión humana, no iniciativa presentada.}

\end{columns}

\end{frame}





% 14
\begin{frame}{Caso II: The Legal Second Brain}
\stagebar{SEGUNDO LABORATORIO}{Superior knowledge coupled with intelligence}
\key{El segundo caso no multiplica agentes. Multiplica conocimiento disponible: un vault jurídico estructurado que permite investigar, conectar evidencia y razonar con procedencia.}
\vspace{0.22cm}
\begin{center}
\begin{tikzpicture}[node distance=.22cm,every node/.style={align=center}]
\node[draw=CambridgeNavy,rounded corners=3pt,fill=Pale,minimum width=2.35cm,minimum height=.92cm,font=\scriptsize\bfseries] (a) {Vault jurídico\\350 documentos};
\node[draw=JudgeTeal,rounded corners=3pt,fill=IdeaPale,minimum width=2.35cm,minimum height=.92cm,font=\scriptsize\bfseries,right=of a] (b) {Cinco grafos\\relacionales};
\node[draw=ExecBlue,rounded corners=3pt,fill=SoftBlue,minimum width=2.35cm,minimum height=.92cm,font=\scriptsize\bfseries,right=of b] (c) {Semillas\\de investigación};
\node[draw=CambridgeGold,rounded corners=3pt,fill=SoftAmber,minimum width=2.35cm,minimum height=.92cm,font=\scriptsize\bfseries,right=of c] (d) {Bucle autónomo\\acotado};
\node[draw=SafeRed,rounded corners=3pt,fill=SoftRed,minimum width=2.35cm,minimum height=.92cm,font=\scriptsize\bfseries,right=of d] (e) {Memorando\\con evidencia};
\foreach \x/\y in {a/b,b/c,c/d,d/e}{\draw[-{Latex},thick,color=Quiet] (\x)--(\y);}
\end{tikzpicture}
\end{center}
\vspace{0.24cm}
\begin{columns}[T,onlytextwidth]
\column{.31\textwidth}\sidecard{JudgeTeal}{NO ES SÓLO BÚSQUEDA}{La pregunta no es qué texto se parece, sino qué evidencia debe investigarse después.}
\column{.31\textwidth}\sidecard{CambridgeGold}{NO ES MEMORIA PASIVA}{El vault conserva documentos, relaciones, procedencia y rutas de investigación.}
\column{.31\textwidth}\sidecard{SafeRed}{NO ES AUTONOMÍA ILIMITADA}{La autonomía es selección adaptativa bajo restricciones explícitas.}
\end{columns}
\end{frame}

% 15
\begin{frame}{Del vault al juicio de investigación}
\stagebar{SECOND BRAIN}{Cuando el conocimiento externo cambia la inteligencia posible}
\begin{columns}[T,onlytextwidth]
\column{.54\textwidth}
\begin{tikzpicture}[node distance=.12cm,every node/.style={align=left}]
\node[draw=CambridgeNavy,fill=Pale,rounded corners=3pt,text width=6.8cm,minimum height=.75cm,inner sep=6pt] (q) {\color{CambridgeNavy}\bfseries\small Consulta jurídica compleja};
\node[draw=JudgeTeal,fill=IdeaPale,rounded corners=3pt,text width=6.8cm,minimum height=.75cm,inner sep=6pt,below=of q] (ss) {\color{JudgeTeal}\bfseries\small Selección de semillas relevantes};
\node[draw=ExecBlue,fill=SoftBlue,rounded corners=3pt,text width=6.8cm,minimum height=.75cm,inner sep=6pt,below=of ss] (g) {\color{ExecBlue}\bfseries\small Expansión por grafos: autoridad, tema, proceso, conflicto, evidencia};
\node[draw=CambridgeGold,fill=SoftAmber,rounded corners=3pt,text width=6.8cm,minimum height=.75cm,inner sep=6pt,below=of g] (u) {\color{CambridgeGold}\bfseries\small Evaluación de incertidumbre y siguiente mejor documento};
\node[draw=SafeRed,fill=SoftRed,rounded corners=3pt,text width=6.8cm,minimum height=.75cm,inner sep=6pt,below=of u] (m) {\color{SafeRed}\bfseries\small Memorando: respuesta + fuentes + límites + ruta recorrida};
\foreach \x/\y in {q/ss,ss/g,g/u,u/m}{\draw[-{Latex},thick,color=Quiet] (\x)--(\y);}
\end{tikzpicture}
\column{.40\textwidth}
\sidecard{JudgeTeal}{CONOCIMIENTO}{Un vault no es almacenamiento: es una arquitectura de relaciones.}
\vspace{0.16cm}
\sidecard{CambridgeGold}{INTELIGENCIA}{El agente decide qué evidencia consultar, no sólo resume lo encontrado.}
\vspace{0.16cm}
\sidecard{SafeRed}{TRAZABILIDAD}{El resultado debe mostrar camino, fuentes, incertidumbres y razones para detenerse.}
\end{columns}
\end{frame}

% 16
\begin{frame}{Dos casos, una lección institucional}
\stagebar{COMPARACIÓN}{Más inteligencia coordinada · más conocimiento utilizable}
\vspace{0.10cm}
\begin{columns}[T,onlytextwidth]
\column{.48\textwidth}
\begin{beamercolorbox}[rounded=true,sep=0.70em,wd=\textwidth]{block body}
{\color{JudgeTeal}\scriptsize\bfseries CASO I · 50 AGENTES}\par\medskip
{\color{CambridgeNavy}\Large\bfseries Superior intelligence}\par\medskip
{\small La capacidad surge de múltiples perspectivas especializadas coordinadas en constelaciones, crítica adversarial y un arnés de gobierno.}\par\medskip
\pill{agentes}\hspace{0.07cm}\pill{roles}\hspace{0.07cm}\pill{coordinación}
\end{beamercolorbox}
\column{.48\textwidth}
\begin{beamercolorbox}[rounded=true,sep=0.70em,wd=\textwidth]{block body}
{\color{CambridgeGold}\scriptsize\bfseries CASO II · SECOND BRAIN}\par\medskip
{\color{CambridgeNavy}\Large\bfseries Superior knowledge + intelligence}\par\medskip
{\small La capacidad surge de un conocimiento externo estructurado que permite explorar relaciones, seleccionar evidencia y construir memorandos auditables.}\par\medskip
\pill{vault}\hspace{0.07cm}\pill{grafos}\hspace{0.07cm}\pill{procedencia}
\end{beamercolorbox}
\end{columns}
\vspace{0.34cm}
\key{La IA jurídica del futuro no será sólo un modelo. Será la combinación gobernada de modelos, memoria, conocimiento, herramientas, agentes y autoridad institucional.}
\end{frame}

\begin{frame}{Qué corresponde hacer ahora}
\stagebar{AGENDA}{Cuatro tareas para pasar de uso individual a capacidad institucional}
\begin{columns}[T,onlytextwidth]
\column{.48\textwidth}
\sidecard{JudgeTeal}{ESTUDIANTES}{Usar IA, pero cultivar criterio: fuente, argumento, institución y responsabilidad.}
\vspace{0.16cm}
\sidecard{ExecBlue}{ESCUELA}{Construir política de herramientas, datos, confidencialidad y experimentación.}
\column{.48\textwidth}
\sidecard{CambridgeGold}{PROFESORES}{Diseñar ejercicios que exijan verificar, comparar, impugnar y defender decisiones.}
\vspace{0.16cm}
\sidecard{SafeRed}{EJERCICIO PROFESIONAL}{Identificar conocimiento, fuentes autorizadas, procesos y decisiones indelegables.}
\end{columns}
\vspace{0.34cm}
\key{El primer paso no es comprar el modelo más sofisticado; es diseñar la arquitectura institucional que permite usarlo responsablemente.}
\end{frame}



% 15

{

\setbeamercolor{background canvas}{bg=CambridgeNavy}

\begin{frame}[plain]

\begin{tikzpicture}[remember picture,overlay]

  \fill[JudgeTeal] (current page.south west) rectangle ([xshift=0.16cm]current page.north west);

  \fill[JudgeTeal!15!CambridgeNavy] ([xshift=-4.35cm]current page.south east) rectangle (current page.north east);

\end{tikzpicture}

\vspace{0.58cm}

{\color{JudgeTeal}\scriptsize\bfseries CIERRE}\par

\vspace{0.74cm}

\begin{minipage}{11.15cm}

{\color{white}\fontsize{17}{20}\selectfont\bfseries

El Derecho lleva siglos construyendo\\[0.04cm]arneses para la inteligencia humana.}

\end{minipage}\par

\vspace{0.48cm}

\begin{minipage}{9.85cm}

{\color{white!86}\small Ahora debe aprender a gobernar otra forma de inteligencia: una que conoce, recomienda, coordina y puede actuar dentro de instituciones.}

\end{minipage}\par

\vspace{0.56cm}

\textcolor{CambridgeGold}{\rule{1.65cm}{0.8pt}}\par

\vspace{0.22cm}

\begin{minipage}{10.35cm}

{\color{CambridgeGold}\fontsize{13}{16}\selectfont\bfseries

En la era de la inteligencia artificial, el Derecho no se vuelve menos importante.}

\end{minipage}\par

\vfill

{\color{white}\fontsize{18}{21}\selectfont\bfseries Se vuelve indispensable.}\par

\vspace{0.34cm}

{\color{white!65}\tiny Alejandro Reynoso}

\vspace{0.38cm}

\end{frame}

}



\end{document}